\par\Needspace{6\baselineskip}\medskip
\hypertarget{native-132-article097}{}\pdfbookmark[2]{97.}{article-097}
\begin{center}97.\end{center}
To illustrate the method developed from Art. 88 onward, we return to the two examples of Art. 87. It scarcely needs pointing out that the meaning assigned so far to the auxiliary angle $\omega$ should be distinguished from the one the same symbol had in Arts. 86 and 87.

I. In the first example we have $f=3^0\,47'\,26''{,}865$ and furthermore
\[
\begin{gathered}
\log\dfrac{r'}r=9{,}991\,4599,\qquad
\log\operatorname{tang}(45^0+\omega)=9{,}997\,8649\,75,\\
\omega=-8'\,27''{,}006.
\end{gathered}
\]
\noindent Hence, by Art. 89,

\begin{center}
\begin{tabular}{@{}ll@{}}
\makebox[17em][l]{$\log\sin\tfrac12 f^2$\hfill .\,.\,.\,.\,.\,.\,.\,.\,.\,.} & $7{,}038\,9972$\\[.3em]
\makebox[17em][l]{$\log\cos f$\hfill .\,.\,.\,.\,.\,.\,.\,.\,.\,.\,.\,.} & $9{,}999\,0488$\\[.3em]
\hline
 & $7{,}039\,9484$\\[.3em]
\multicolumn{2}{l}{$=\log0{,}001\,0963\,480$}\\[.3em]
\end{tabular}
\end{center}

\begin{center}
\begin{tabular}{@{}ll@{}}
\makebox[17em][l]{$\log\operatorname{tang}2\omega^2$\hfill .\,.\,.\,.\,.\,.\,.\,.} & $5{,}383\,2428$\\[.3em]
\makebox[17em][l]{$\log\cos f$\hfill .\,.\,.\,.\,.\,.\,.\,.\,.\,.\,.} & $9{,}999\,0488$\\[.3em]
\hline
 & $5{,}384\,1940$\\[.3em]
\multicolumn{2}{l}{$=\log0{,}000\,0242\,211$}\\[.3em]
\end{tabular}
\end{center}

\noindent\SourceLeaf{133}and consequently $l=0{,}001\,1205\,691$, $\tfrac56+l=0{,}834\,4539$. Further, we obtain $\log kt=9{,}576\,6974$

\begin{center}
\begin{tabular}{@{}ll@{}}
\makebox[17em][l]{$2\log kt$\hfill .\,.\,.\,.\,.\,.\,.\,.\,.\,.\,.\,.\,.\,.} & $9{,}153\,3948$\\[.3em]
\makebox[17em][l]{$\mathrm{C.}\tfrac32\log rr'$\hfill .\,.\,.\,.\,.\,.\,.\,.\,.\,.\,.} & $9{,}020\,5181$\\[.3em]
\makebox[17em][l]{$\mathrm{C.}\log8\cos f^3$\hfill .\,.\,.\,.\,.\,.\,.\,.\,.} & $9{,}099\,7636$\\[.3em]
\hline
\makebox[17em][l]{$\log mm$\hfill .\,.\,.\,.\,.\,.\,.\,.\,.\,.\,.\,.\,.\,.} & $7{,}273\,6765$\\[.3em]
\makebox[17em][l]{$\log(\tfrac56+l)$\hfill .\,.\,.\,.\,.\,.\,.\,.\,.\,.\,.\,.} & $9{,}921\,4023$\\[.3em]
\hline
 & $7{,}352\,2742$\\[.3em]
\end{tabular}
\end{center}

\noindent The approximate value is therefore $h=0{,}002\,2504\,7$, to which $\log yy=0{,}002\,1633$ corresponds in our Table II. Thus we have $\log\dfrac{mm}{yy}=7{,}271\,5132$, or $\dfrac{mm}{yy}=0{,}001\,8685\,87$, which by formula 16 gives $x=0{,}000\,7480\,179$. Since Table III shows $\xi$ to be entirely negligible, the values found for $h$, $y$, and $x$ require no correction. We then determine the elements as follows.

\begin{center}
\begin{tabular}{@{}ll@{}}
\makebox[17em][l]{$\log x$\hfill .\,.\,.\,.\,.\,.\,.\,.\,.\,.\,.\,.\,.\,.\,.\,.} & $6{,}873\,9120$\\[.3em]
\makebox[17em][l]{$\log\sin\tfrac12 g$\hfill .\,.\,.\,.\,.\,.\,.\,.\,.\,.\,.\,.} & $8{,}436\,9560$\\[.3em]
\end{tabular}
\end{center}

\[
\begin{gathered}
\tfrac12 g=1^0\,34'\,2''{,}0286,\qquad
\tfrac12(f+g)=3^0\,27'\,45''{,}4611,\\
\tfrac12(f-g)=19'\,41''{,}4039.
\end{gathered}
\]
\Needspace{19\baselineskip}
\noindent Hence formulas 27, 28, 29, and 30 give

\begin{center}
\begin{tabular}{@{}ll@{}}
\makebox[17em][l]{$\log\operatorname{tang}2\omega$\hfill .\,.\,.\,.\,.\,.\,.\,.\,.\,.\,.} & $7{,}691\,6214_n$\\[.3em]
\makebox[17em][l]{$\log\cos\tfrac12(f+g)$\hfill .\,.\,.\,.\,.\,.\,.\,.} & $9{,}999\,2065$\\[.3em]
\makebox[17em][l]{$\log\cos\tfrac12(f-g)$\hfill .\,.\,.\,.\,.\,.\,.\,.} & $9{,}999\,9929$\\[.3em]
\hline\noalign{\vskip .3em}
\makebox[17em][l]{$\log P\sin\tfrac12(F-G)$\hfill .\,.\,.\,.\,.} & $7{,}690\,8279_n$\\[.3em]
\makebox[17em][l]{$\log P\cos\tfrac12(F-G)$\hfill .\,.\,.\,.\,.} & $8{,}781\,0240$\\[.3em]
\hline\noalign{\vskip .3em}
$\tfrac12(F-G)=$ & $-4^0\,38'\,41''{,}54$\\[.3em]
$\tfrac12(F+G)=$ & $319\quad21\quad38{,}05$\\[.3em]
\hline
$F=$ & $314\quad42\quad56{,}51$\\[.3em]
$v=$ & $310\quad55\quad29{,}64$\\[.3em]
$v'=$ & $318\quad30\quad23{,}37$\\[.3em]
$G=$ & $324\quad0\quad19{,}59$\\[.3em]
$E=$ & $320\quad52\quad15{,}53$\\[.3em]
$E'=$ & $327\quad8\quad23{,}65$\\[.3em]
\end{tabular}
\end{center}

\begin{center}
\begin{tabular}{@{}ll@{}}
\makebox[17em][l]{$\mathrm{C.}\log\cos2\omega$\hfill .\,.\,.\,.\,.\,.\,.\,.\,.\,.} & $0{,}000\,0052$\\[.3em]
\makebox[17em][l]{$\log\sin\tfrac12(f+g)$\hfill .\,.\,.\,.\,.\,.\,.\,.} & $8{,}781\,0188$\\[.3em]
\makebox[17em][l]{$\log\sin\tfrac12(f-g)$\hfill .\,.\,.\,.\,.\,.\,.\,.} & $7{,}757\,9709$\\[.3em]
\hline\noalign{\vskip .3em}
\makebox[17em][l]{$\log Q\sin\tfrac12(F+G)$\hfill .\,.\,.\,.\,.} & $7{,}691\,6143_n$\\[.3em]
\makebox[17em][l]{$\log Q\cos\tfrac12(F+G)$\hfill .\,.\,.\,.\,.} & $7{,}757\,9761$\\[.3em]
\hline\noalign{\vskip .3em}
\makebox[17em][l]{$\log P=\log R\cos\tfrac12\varphi$\hfill .\,.\,.} & $8{,}782\,4527$\\[.3em]
\makebox[17em][l]{$\log Q=\log R\sin\tfrac12\varphi$\hfill .\,.\,.} & $7{,}877\,8355$\\[.3em]
\hline\noalign{\vskip .3em}
$\text{Hence }\tfrac12\varphi=$ & $7^0\,6'\,0''{,}935$\\[.3em]
$\varphi=$ & $14\quad12\quad1{,}87$\\[.3em]
\makebox[17em][l]{$\log R$\hfill .\,.\,.\,.\,.\,.\,.\,.\,.\,.\,.\,.\,.\,.\,.} & $8{,}785\,7960$\\[.3em]
\end{tabular}
\end{center}

\Needspace{9\baselineskip}
\begin{center}\small To check the calculation\end{center}

\begin{center}
\begin{tabular}{@{}ll@{}}
\makebox[17em][l]{$\tfrac12\log2\cos f$\hfill .\,.\,.\,.\,.\,.\,.\,.\,.\,.\,.} & $0{,}150\,0394$\\[.3em]
\makebox[17em][l]{$\tfrac12\log(l+x)=\log\dfrac m y$\hfill .\,.\,.\,.} & $8{,}635\,7566$\\[.3em]
\noalign{\vskip .3em}\hline
 & $8{,}785\,7960$\\[.3em]
\end{tabular}
\end{center}

\SourceLeaf{134}

\begin{center}
\begin{tabular}{@{}ll@{}}
\makebox[17em][l]{$\tfrac12\log rr'$\hfill .\,.\,.\,.\,.\,.\,.\,.\,.\,.\,.\,.\,.\,.} & $0{,}326\,4939$\\[.3em]
\makebox[17em][l]{$\log\sin f$\hfill .\,.\,.\,.\,.\,.\,.\,.\,.\,.\,.\,.\,.\,.\,.} & $8{,}820\,2909$\\[.3em]
\makebox[17em][l]{$\mathrm{C.}\log\sin g$\hfill .\,.\,.\,.\,.\,.\,.\,.\,.\,.\,.\,.} & $1{,}262\,1765$\\[.3em]
\hline
\makebox[17em][l]{$\log b$\hfill .\,.\,.\,.\,.\,.\,.\,.\,.\,.\,.\,.\,.\,.\,.\,.\,.} & $0{,}408\,9613$\\[.3em]
\makebox[17em][l]{$\log\cos\varphi$\hfill .\,.\,.\,.\,.\,.\,.\,.\,.\,.\,.\,.\,.\,.} & $9{,}986\,5224$\\[.3em]
\hline
\makebox[17em][l]{$\log p$\hfill .\,.\,.\,.\,.\,.\,.\,.\,.\,.\,.\,.\,.\,.\,.\,.\,.} & $0{,}395\,4837$\\[.3em]
\makebox[17em][l]{$\log a$\hfill .\,.\,.\,.\,.\,.\,.\,.\,.\,.\,.\,.\,.\,.\,.\,.\,.} & $0{,}422\,4389$\\[.3em]
\end{tabular}
\end{center}

\begin{center}
\begin{tabular}{@{}ll@{}}
\makebox[17em][l]{$\log k$\hfill .\,.\,.\,.\,.\,.\,.\,.\,.\,.\,.\,.\,.\,.\,.\,.\,.} & $3{,}550\,0066$\\[.3em]
\makebox[17em][l]{$\tfrac32\log a$\hfill .\,.\,.\,.\,.\,.\,.\,.\,.\,.\,.\,.\,.\,.\,.\,.} & $0{,}633\,6584$\\[.3em]
\hline
 & $2{,}916\,3482$\\[.3em]
\makebox[17em][l]{$\log t$\hfill .\,.\,.\,.\,.\,.\,.\,.\,.\,.\,.\,.\,.\,.\,.\,.\,.} & $1{,}341\,1160$\\[.3em]
\hline
 & $4{,}257\,4642$\\[.3em]
\end{tabular}
\end{center}

\noindent The mean daily motion is therefore $=824''{,}7990$. The mean motion in the time $t=18091''{,}07=5^0\,1'\,31''{,}07$.

\begin{center}
\begin{tabular}{@{}ll@{}}
\makebox[17em][l]{$\log\sin\varphi$\hfill .\,.\,.\,.\,.\,.\,.\,.\,.\,.\,.\,.\,.\,.} & $9{,}389\,7262$\\[.3em]
\makebox[17em][l]{$\log206265$\hfill .\,.\,.\,.\,.\,.\,.\,.\,.\,.\,.\,.} & $5{,}314\,4251$\\[.3em]
\hline
\makebox[17em][l]{$\log e\text{ in arcseconds}$\hfill .\,.\,.\,.\,.\,.\,.\,.} & $4{,}704\,1513$\\[.3em]
\makebox[17em][l]{$\mathop{\mathrm{lo}\bar{\mathrm{g}}}\nolimits\sin E$\hfill .\,.\,.\,.\,.\,.\,.\,.\,.\,.\,.\,.\,.\,.} & $9{,}800\,0767_n$\\[.3em]
\makebox[17em][l]{$\log\sin E'$\hfill .\,.\,.\,.\,.\,.\,.\,.\,.\,.\,.\,.\,.} & $9{,}734\,4714_n$\\[.3em]
\hline
\makebox[17em][l]{$\log e\sin E$\hfill .\,.\,.\,.\,.\,.\,.\,.\,.\,.\,.\,.\,.} & $4{,}504\,2280_n$\\[.3em]
\makebox[17em][l]{$\log e\sin E'$\hfill .\,.\,.\,.\,.\,.\,.\,.\,.\,.\,.\,.} & $4{,}438\,6227_n$\\[.3em]
\end{tabular}
\end{center}

\[
\begin{aligned}
e\sin E&=-31932''{,}14=-8^0\,52'\,12''{,}14\\
e\sin E'&=-27455{,}08=-7\quad37\quad35{,}08
\end{aligned}
\]
\Needspace{8\baselineskip}
\begin{center}Hence the mean anomaly\end{center}

\begin{center}
\begin{tabular}{@{}ll@{}}
$\text{for the first position}=$ & $329^0\,44'\,27''{,}67$\\[.3em]
$\text{for the second}=$ & $334\quad45\quad58{,}73$\\[.3em]
\hline
$\text{Difference}=$ & $5\quad1\quad31{,}06$\\[.3em]
\end{tabular}
\end{center}

II. In the second example we obtain $f=31^0\,27'\,38''{,}32$, $\omega=-21'\,50''{,}565$, $l=0{,}086\,3565\,9$, $\log mm=9{,}353\,0651$, and the quotient $\dfrac{mm}{\frac56+l}$ gives the approximate value $h=0{,}245\,1454$. To this corresponds $\log yy=0{,}172\,2663$ in Table II, from which we deduce $\dfrac{mm}{yy}=0{,}151\,6347\,7$ and $x=0{,}065\,2781\,8$. The latter yields $\xi=0{,}000\,2531$ from Table III. Using this value gives the corrected values $h=0{,}245\,0779$, $\log yy=0{,}172\,2303$, $\dfrac{mm}{yy}=0{,}151\,6473\,5$, $x=0{,}065\,2907\,6$, and $\xi=0{,}000\,2532$. If we repeated the calculation with this value of $\xi$, which differs from the previous one by just one unit in the seventh decimal place, $h$, $\log yy$, and $x$ would show no appreciable change. The value found for $x$ is therefore already the true one, and we may proceed at once to determine the elements. We shall pass over that step here, since it proceeds exactly as in the preceding example.

III. It will also be useful to illustrate the other case, in which $\cos f$ is negative, with an example. Let $v'-v=224^0\,0'\,0''$, or $f=112^0\,0'\,0''$, and let $\log r=0{,}139\,4892$, $\log r'=0{,}397\,8794$, and $t=206{,}80919$ days. Here we find $\omega=+4^0\,14'\,43''{,}82$, $L=1{,}894\,2298$, $\log MM=0{,}672\,4333$, and as a first approximation $\log H=0{,}646\,7603$. Solving equation $15^*$ then gives $Y=1{,}591\,432$, and hence $x=0{,}037\,037$, to which $\xi=0{,}000\,0801$ corresponds in Table III. \SourceLeaf{135}This gives the corrected values $\log H=0{,}646\,7931$, $Y=1{,}591\,5107$, $x=0{,}037\,2195$, and $\xi=0{,}000\,0809$. Repeating the calculation with this value of $\xi$ gives $x=0{,}037\,2213$. This value requires no further correction, since it leaves $\xi$ unchanged. We then find $\tfrac12 g=11^0\,7'\,25''{,}40$ and hence, just as in example I,

\begin{center}
\begin{tabular}{@{}ll@{}}
$\tfrac12(F-G)=$ & $3^0\,33'\,53''{,}62$\\[.3em]
$\tfrac12(F+G)=$ & $8\quad26\quad6{,}38$\\[.3em]
$F=$ & $12\quad0\quad0{,}00$\\[.3em]
$v=$ & $-100\quad0\quad0{,}00$\\[.3em]
$v'=$ & $+124\quad0\quad0{,}00$\\[.3em]
$G=$ & $4\quad52\quad12{,}76$\\[.3em]
$E=$ & $-17\quad22\quad38{,}04$\\[.3em]
$E'=$ & $+27\quad7\quad3{,}56$\\[.3em]
\end{tabular}
\end{center}

\begin{center}
\begin{tabular}{@{}ll@{}}
\makebox[17em][l]{$\log P=\log R\cos\tfrac12\varphi$\hfill .\,.\,.\,.} & $9{,}970\,0508$\\[.3em]
\makebox[17em][l]{$\log Q=\log R\sin\tfrac12\varphi$\hfill .\,.\,.\,.} & $9{,}858\,0553$\\[.3em]
$\tfrac12\varphi=$ & $37^0\,41'\,34''{,}27$\\[.3em]
$\varphi=$ & $75\quad23\quad8{,}54$\\[.3em]
\makebox[17em][l]{$\log R$\hfill .\,.\,.\,.\,.\,.\,.\,.\,.\,.\,.\,.\,.\,.\,.\,.\,.} & $0{,}071\,7097$\\[.3em]
\end{tabular}
\end{center}

\Needspace{6\baselineskip}
\begin{center}\small As a check on the calculation we obtain\end{center}

\begin{center}
\begin{tabular}{@{}ll@{}}
\makebox[17em][l]{$\log\dfrac M Y\BareRoot{(-2\cos f)}$\hfill .\,.\,.\,.\,.\,.\,.} & $0{,}071\,7097$\\[.3em]
\end{tabular}
\end{center}

\noindent In orbits of such high eccentricity the angle $\varphi$ can be computed a little more accurately by formula $19^*$, which in our example gives $\varphi=75^0\,23'\,8''{,}57$. The eccentricity $e$ is also determined more precisely by the formula $1-2\sin(45^0-\tfrac12\varphi)^2$ than by $\sin\varphi$. The former gives $e=0{,}967\,6463\,0$.

Further, formula 1 gives $\log b=0{,}657\,6611$, whence $\log p=0{,}0595967$, $\log a=1{,}255\,7255$, and for the logarithm of the distance at perihelion we obtain
\[
\begin{aligned}
{}={}&\log\dfrac p{1+e}=\log(a(1-e))\\
{}={}&\log(b\operatorname{tang}(45^0-\tfrac12\varphi))=9{,}765\,6496.
\end{aligned}
\]

For orbits so closely resembling a parabola, it is usual to specify the time of passage through perihelion in place of the epoch of the mean anomaly. Once the elements are known, the intervals between this time and the times corresponding to the two given positions can be determined by the method given in Art. 41. Their difference or sum, according as perihelion lies outside or between the two given positions, serves to check the calculation, since it must agree with the time $t$. — The numbers in this third example were based on the elements assumed in the example of Arts. 38 and 43. Indeed, that very example supplied our first position. The slight differences in the elements obtained here arose solely from the limited precision of the logarithmic and trigonometric tables.
